% !PRINTED 147
\weberpage{147}
Wäre nämlich $Q_0$ in diesen $n-1$ Ziffern intransitiv, so müsste nach $\alpha$)

\[
\begin{aligned}
&\text{entweder } \Pi(n-2) \text{ theilbar durch } g,\\
&\text{oder } \Pi(n-3)\,\Pi(2) \geqq g
\end{aligned}
\]

sein; also nach (9):

\[
\begin{aligned}
&\text{entweder } \frac{n}{n-1} \text{ eine ganze Zahl, also } n \geqq 2(n-1)\\
&\text{oder } n^2 - 5n + 2 \leqq 0.
\end{aligned}
\]

Das Eine ist unmöglich, wenn $n > 2$ ist, das Andere, wenn $n > 4$ ist.

Wir sehen von dem Falle $n = 4$ ab und betrachten jetzt den Theiler $Q_{0,1}$ von
$Q$, der die beiden Ziffern 0, 1 ungeändert lässt.

Da $Q_{0,1}$, als Permutationsgruppe der $n-1$ Ziffern betrachtet, ein transitiver
Theiler von $Q_0$ ist, so ist, wie man durch nochmalige Anwendung der Zerlegung (8)
schliesst, der Grad $g_1$ von $Q_{0,1}$ gleich $g : n-1$, also

\weq{10}{g_1 = \frac{\Pi(n-1)}{n(n-1)} = \frac{\Pi(n-2)}{n}.}

Daraus schliessen wir ähnlich wie oben:

\gesperrt{$\delta$) Ist $n$ nicht $= 6$, so sind die $n-2$ Ziffern
$2, 3, \ldots, n-1$ durch $Q_{0,1}$ transitiv verbunden.}

Denn wären sie es nicht, so müsste nach $\alpha$):

\[
\begin{aligned}
&\text{entweder } \Pi(n-3) \text{ theilbar durch } g_1,\\
&\text{oder } \Pi(n-4)\,\Pi(2) \geqq g_1
\end{aligned}
\]

sein; also nach (10):

\[
\begin{aligned}
&\text{entweder } \frac{n}{n-2} \text{ eine ganze Zahl, also } n \geqq 2n - 4,\\
&\text{oder } n^2 - 7n + 6 \leqq 0.
\end{aligned}
\]

Das Erste ist nicht möglich, wenn $n > 4$, das Zweite, wenn $n > 6$ ist.

Ist nun $n > 6$, so gilt noch Folgendes:

\gesperrt{$\varepsilon$) Die Gruppe $Q_{0,1,2}$, die die drei Ziffern 0, 1, 2
ungeändert lässt, hat den Grad}

\[ g_2 = \frac{g_1}{n-2} = \frac{\Pi(n-3)}{n}, \]

\gesperrt{und es ist nicht möglich, dass durch die ganze Gruppe $Q_{0,1,2}$ noch eine
vierte Ziffer 3 ungeändert bleibt.}
